%%%%%%%%%%%%%%%%%%%%%%%%%%%%%%%%%%%%%%%%%
% Beamer Presentation
% LaTeX Template
% Version 1.0 (10/11/12)
%
% This template has been downloaded from:
% http://www.LaTeXTemplates.com
%
% License:
% CC BY-NC-SA 3.0 (http://creativecommons.org/licenses/by-nc-sa/3.0/)
%
%%%%%%%%%%%%%%%%%%%%%%%%%%%%%%%%%%%%%%%%%

%----------------------------------------------------------------------------------------
%	PACKAGES AND THEMES
%----------------------------------------------------------------------------------------

\documentclass[]{beamer}

\mode<presentation> {
	
	% The Beamer class comes with a number of default slide themes
	% which change the colors and layouts of slides. Below this is a list
	% of all the themes, uncomment each in turn to see what they look like.
	
	%\usetheme{default}
	%\usetheme{AnnArbor}
	%\usetheme{Antibes}
	%\usetheme{Bergen}
	%\usetheme{Berkeley}
	%\usetheme{Berlin}
	%\usetheme{Boadilla}
	%\usetheme{CambridgeUS}
	%\usetheme{Copenhagen}
	%\usetheme{Darmstadt}
	%\usetheme{Dresden}
	%\usetheme{Frankfurt}
	%\usetheme{Goettingen}
	%\usetheme{Hannover}
	%\usetheme{Ilmenau}
	%\usetheme{JuanLesPins}
	%\usetheme{Luebeck}
	\usetheme{Madrid}
	%\usetheme{Malmoe}
	%\usetheme{Marburg}
	%\usetheme{Montpellier}
	%\usetheme{PaloAlto}
	%\usetheme{Pittsburgh}
	%\usetheme{Rochester}
	%\usetheme{Singapore}
	%\usetheme{Szeged}
	%\usetheme{Warsaw}
	
	% As well as themes, the Beamer class has a number of color themes
	% for any slide theme. Uncomment each of these in turn to see how it
	% changes the colors of your current slide theme.
	
	%\usecolortheme{albatross}
	%\usecolortheme{beaver}
	%\usecolortheme{beetle}
	%\usecolortheme{crane}
	%\usecolortheme{dolphin}
	%\usecolortheme{dove}
	%\usecolortheme{fly}
	%\usecolortheme{lily}
	%\usecolortheme{orchid}
	%\usecolortheme{rose}
	%\usecolortheme{seagull}
	%\usecolortheme{seahorse}
	%\usecolortheme{whale}
	%\usecolortheme{wolverine}
	
	%\setbeamertemplate{footline} % To remove the footer line in all slides uncomment this line
	%\setbeamertemplate{footline}[page number] % To replace the footer line in all slides with a simple slide count uncomment this line
	
	\setbeamertemplate{navigation symbols}{} % To remove the navigation symbols from the bottom of all slides uncomment this line
}

\makeatletter
\setbeamertemplate{footline}
{
	\leavevmode%
	\hbox{%
		\begin{beamercolorbox}[wd=.333333\paperwidth,ht=2.25ex,dp=1ex,center]{author in head/foot}%
			\usebeamerfont{author in head/foot}\insertsection
		\end{beamercolorbox}%
		\begin{beamercolorbox}[wd=.333333\paperwidth,ht=2.25ex,dp=1ex,center]{title in head/foot}%
			\usebeamerfont{title in head/foot}\insertsubsection
		\end{beamercolorbox}%
		\begin{beamercolorbox}[wd=.333333\paperwidth,ht=2.25ex,dp=1ex,right]{date in head/foot}%
			\usebeamerfont{date in head/foot}\insertshortdate{}\hspace*{2em}
			\insertframenumber{} / \inserttotalframenumber\hspace*{2ex} 
	\end{beamercolorbox}}%
	\vskip0pt%
}
\makeatother

\usepackage{graphicx} % Allows including images
\usepackage{booktabs} % Allows the use of \toprule, \midrule and \bottomrule in tables
\usepackage{amssymb,amsmath,dsfont,physics}
\usepackage[english]{babel}
\usefonttheme[onlymath]{serif}
\usepackage{multicol}
\usepackage{datetime,xcolor}

% Definitions
\newcommand{\C}[1]{\ensuremath{\mathbb{C}^{#1}}}
\newcommand{\R}[1]{\ensuremath{\mathbb{R}^{#1}}}
\newcommand{\Exp}[1]{\ensuremath{E\qty[#1]}}
\newcommand{\Cov}[1]{\ensuremath{\text{Cov}\qty(#1)}}
\newcommand{\myspace}{\vspace{3mm}}
\newcommand{\herm}[1]{\ensuremath{#1^\dagger}}
\newcommand{\cscg}{circularly symmetric Complex Gaussian}

% Custom colors
\definecolor{DEIred}{RGB}{153,14,32}
\definecolor{darkDEI1}{RGB}{117,11,24}
\definecolor{darkDEI2}{RGB}{78,7,16}
\definecolor{darkDEI3}{RGB}{117,11,24}

\setbeamercolor{palette primary}{bg=DEIred,fg=white}
\setbeamercolor{palette secondary}{bg=darkDEI1,fg=white}
\setbeamercolor{palette tertiary}{bg=darkDEI2,fg=white}
\setbeamercolor{palette quaternary}{bg=darkDEI3,fg=white}
\setbeamercolor{structure}{fg=DEIred} % itemize, enumerate, etc
\setbeamercolor{section in toc}{fg=black} % TOC sections
% Override palette coloring with secondary
\setbeamercolor{subsection in head/foot}{bg=blue,fg=white}

%----------------------------------------------------------------------------------------
%	TITLE PAGE
%----------------------------------------------------------------------------------------

\title[Notation]{Capacity of Multi-antenna Gaussian Channels\\
	(I. E. Telatar, 1999)} % The short title appears at the bottom of every slide, the full title is only on the title page

\author[Lecci]{Mattia Lecci\inst{1}\inst{2}} % Your name

\institute[UniPD/UPC] % Your institution as it will appear on the bottom of every slide, may be shorthand to save space
{
	\inst{1}
	Universit\`a degli Studi di Padova\\
	Department of Information Engineering\\
	(DEI-UniPD)
	\and
	\inst{2}
	Universitat Polit\`{e}cnica de Catalunya\\ % Your institution for the title page
	Escola T\`ecnica Superior d'Engenyeria de Telecomunicaci\'o  de Barcelona\\
	(UPC-ETSETB)
	%\medskip
	%\textit{javier.fonollosa@upc.edu} % Your email address
}
\date{January 2018} % Date, can be changed to a custom date


%------------------------------------------------------------
%------------------------------------------------------------
%The next block of commands puts the table of contents at the 
%beginning of each section and highlights the current section:

\AtBeginSection[]
{
	\begin{frame}
	\frametitle{Table of Contents}
	\tableofcontents[currentsection]
\end{frame}
}
%------------------------------------------------------------


\begin{document}
\begin{frame}
\titlepage % Print the title page as the first slide
\begin{figure}
	\hspace{2cm}
	\includegraphics[height = 1.2cm]{img/dei_logo}
	\hfill
	\includegraphics[height = 1.2cm]{img/ETSETB-positiu-p3005}
	\hspace{1cm}
\end{figure}
\end{frame}

\begin{frame}
\frametitle{Overview} % Table of contents slide, comment this block out to remove it
%\begin{multicols}{2}
	\tableofcontents
%\end{multicols}
\end{frame}

\section{Introduction}

\subsection{Abstract}
%%%%%%%%%%%%%%%%%%%%%%%%%%%%%%%%%%%%%%%%%%%%%%%%%%%%%%%%%
\begin{frame}{Abstract}
In this presentation I will talk about the capacity of a single-user Gaussian channel with multiple receiving and/or transmitting antennas \cite{telatar99} (also known as \alert{MIMO} channel).

\myspace
I will talk about 3 cases:
\begin{itemize}
	\item Deterministic channel
	\item Random i.i.d. ergodic channel
	\item \textit{Bonus}: multi-user case
\end{itemize}

\end{frame}

%%%%%%%%%%%%%%%%%%%%%%%%%%%%%%%%%%%%%%%%%%%%%%%%%%%%%%%%%%
\subsection{Notation and Assumptions}
\begin{frame}[allowframebreaks]{Notation}
The notation adopted is halfway between the one used during the course and the one from the original paper.

\myspace
Let's denote with $t$ the number of transmitting and with $r$ the number of receiving antennas.

\myspace
We will consider the classical linear model, where $\vb{x}\in \C{t}$ is the transmitted vector and $\vb{y}\in \C{r}$ is the received vector, $H\in \C{r\times t}$ is the complex channel matrix and $\vb{n}\in \C{r}$ is the noise
$$\vb{y} = H\vb{x} + \vb{n}$$

%%%%%%%%%%%%%%%%%%%%%%%%%%%%%%%%%%%%%%%%%%%%%%%%%%%%%%%%%%
\framebreak

We assume noise at different receivers to be independent and normalized, i.e. $\Exp{\vb{n} \herm{\vb{n}}} = I_r$.\\
The dag (\dag) notation is used for the conjugate-transpose operation.

\myspace
We have the power constraint
$$\Exp{\herm{\vb{x}}\vb{x}} = \Exp{\tr(\vb{x}\herm{\vb{x}})} = \tr(\Exp{\vb{x}\herm{\vb{x}}})\leq P$$

\myspace
A complex random vector (r.ve.) $\vb{x}\in \C{n}$ is said to be Complex Gaussian if its real extension $\vh{x} \triangleq \qty[\begin{array}{c}
\Re{\vb{x}}\\
\Im{\vb{x}}
\end{array}]
\in \R{2n}$ is Gaussian.

\myspace
The r.ve. $\vb{x}$ will have \textit{mean} and \textit{covariance} respectively $\mu= \Exp{\vb{x}}$ and $Q = \Cov{Q} = \Exp{(\vb{x}-\mu)\herm{(\vb{x}-\mu)}}$

%%%%%%%%%%%%%%%%%%%%%%%%%%%%%%%%%%%%%%%%%%%%%%%%%%%%%%%%%%
\framebreak

Defining for a complex matrix $A$
$$\hat{A} = \qty[\begin{array}{lr}
\Re(A) & -\Im(A)\\
\Im(A) & \Re(A)
\end{array}]$$
we say that a Complex Gaussian r.ve. is \textit{circularly symmetric} if
$$\Cov{\vh{x}} = \Exp{(\vh{x}-\hat{\mu})(\vh{x}-\hat{\mu})^T} = \frac{1}{2} \hat{Q}$$

\begin{block}{pdf of \cscg}
	The pdf of a \cscg is
	\begin{align*}
	\begin{split}
	\gamma_{\mu,Q} &= \det(\pi\hat{Q})^{-\frac{1}{2}} e^{-(\hat{x}-\hat{\mu})^T\hat{Q}(\hat{x}-\hat{\mu})}\\
	&= \det(\pi Q)^{-1} e^{-\herm{(x-\mu)}Q(x-\mu)}
	\end{split}
	\end{align*}
\end{block}

\end{frame}

%%%%%%%%%%%%%%%%%%%%%%%%%%%%%%%%%%%%%%%%%%%%%%%%%%%%%%%%%%
\subsection{Properties and Lemmas}
\begin{frame}[allowframebreaks]{Properties and Lemmas}
\begin{block}{Lemma 1}
	The following properties hold:
	\begin{subequations}
	\begin{gather}
	C=AB \iff \hat{C} = \hat{A}\hat{B}\\
	C = A+B \iff \hat{C} = \hat{A} + \hat{B}\\
	C = \herm{A} \iff \hat{C}=\hat{A}^T\\
	C=A^{-1} \iff \hat{C} = \hat{A}^{-1}\\
	\det(\hat{A}) = |\det(A)|^2 = \det(A \herm{A})\\
	z = x+y \iff \hat{z} = \hat{x}+\hat{y}\\
	y=Ax \iff \hat{y} = \hat{A}\hat{x}\\
	\Re{\herm{x}y} = \hat{x}^T\hat{y}
	\end{gather}
	\end{subequations}
\end{block}

%%%%%%%%%%%%%%%%%%%%%%%%%%%%%%%%%%%%%%%%%%%%%%%%%%%%%%%%%%
\framebreak

\begin{block}{Corollary 1}
A matrix $U\in \C{n\times n}$ is unitary if and only if $\hat{U}\in \R{2n\times 2n}$ is orthonormal.
\end{block}
\begin{block}{Corollary 2}
If $Q\in \C{n\times n}$ is positive semi-definite, then so is $\hat{Q}\in \R{2n\times 2n}$.
\end{block}
\begin{alertblock}{Lemma 2}\justify
Suppose the complex r.ve. $\vb{x}\in \C{n}$ is
zero-mean and $\Cov{\vb{x}}=Q$.\\
Then the differential entropy of $\vb{x}$ satisfies $h(\vb{x}) \leq \log\det(\pi e Q)$ with equality if and only if $\vb{x}$ is \cscg.
\end{alertblock}
In other words, \cscg r.ve. are entropy maximizers for the class of complex random vectors.

%%%%%%%%%%%%%%%%%%%%%%%%%%%%%%%%%%%%%%%%%%%%%%%%%%%%%%%%%%
\framebreak

\begin{block}{Lemma 3}
lf $\vb{x}\in \C{n}$ is a \cscg then so is $\vb{y} = A\vb{x}$ for any $A \in \C{m\times n}$.
\end{block}
\begin{block}{Lemma 4}\label{lemma:4}
	f $\vb{x}$ and $\vb{y}$ are independent {\cscg s}, then $\vb{z}=\vb{x}+\vb{y}$ is also \cscg.
\end{block}

\end{frame}

%%%%%%%%%%%%%%%%%%%%%%%%%%%%%%%%%%%%%%%%%%%%%%%%%%%%%%%%%%
%%%%%%%%%%%%%%%%%%%%%%%%%%%%%%%%%%%%%%%%%%%%%%%%%%%%%%%%%%

\begin{frame}{References}
\nocite{*}
\bibliographystyle{ieeetr}
\bibliography{mybib}
\end{frame}

\end{document} 